\hspace{3mm}

\centerline{\bf \Huge Квадратичные иррациональности}

\begin{flushright}
    \scriptsize{}
\end{flushright}

\textbf{Определение.} \textit{Квадратичной иррациональностью} называется иррациональное число, которое можно представить в виде $a+b\sqrt{n}$, где $a$ и $b$ $-$ рациональные числа, а $n$ $-$ натуральное, свободное от квадратов.

\textbf{Другое определение.} \textit{Квадратичной иррациональностью} называется иррациональное число, являющееся корнем квадратного уравнения с рациональными коэффициентами.

\textbf{Определение.} Число $\overline{x} = a - b\sqrt{n}$ называется сопряженным к квадратичной иррациональности $ x = a + b\sqrt{n}.$ Обозначается с чёрточкой наверху.

\problem{1}
Любая квадратичная иррациональность единственным образом представляется в виде $a+b\sqrt{n}$. Иначе говоря, если две квадратичные иррациональности $a+b\sqrt{n}$ и $c+d\sqrt{m}$ равны это означает, что $a=c,\ b=d,\ m=n.$

\problem{2}
\textit{Лемма о сопряжённом корне.} Если число $x = a+b\sqrt{m}$ является корнем многочлена с рациональными коэффициентами, то тогда $\overline{x}$ тоже является корнем.

\problem{3}
Докажите, что $(\overline{a+b\sqrt{n}})(\overline{c+d\sqrt{n}}) = \overline{(a+b\sqrt{n})(c+d\sqrt{n})}$

\problem{4}
Найдите у чисел 

а)$(6+\sqrt{35})^{1999};$ 

б)$(6 + \sqrt{37})^{1999};$

в)$(6 + \sqrt{37})^{2000}$

первые 1000 знаков после запятой.

\problem{5}
Докажите, что уравнение   $(x+y)^4 + (z+t)^4 = 2 + \sqrt{5}$   не имеет решений в рациональных числах.

\problem{6} 
Пусть $p, q, r$ $-$ различные простые числа. Докажите, что $\sqrt{p}, \sqrt{q}, \sqrt{r}$ не могут быть последовательными членами одной арифметические прогрессии.